\begin{proposition}[I4: adjacent arrows in every dimension; Proved]\label{i:adjacent}
Let $\mathfrak p\subset\mathfrak q$ be an immediate specialization
in a regular noetherian local ring $A$. Put $R=A_{\mathfrak q}$,
$\mathfrak m=\mathfrak qR$, and $T=\widehat{A_{\mathfrak p}}$.
For the residual differential $\delta_{\mathfrak p\mathfrak q}$,
the map
\[
 T\longrightarrow\Hom_A(E_{\mathfrak p},E_{\mathfrak q}),
 \qquad c\longmapsto\delta_{\mathfrak p\mathfrak q}c
\]
is injective and dense. Its induced topology has basis
\[
 \mathfrak m^N R+(\mathfrak pR_{\mathfrak p})^n T,
 \qquad N,n\geq1.
\]
The Hom module is its separated completion for this topology.
\end{proposition}
\begin{proof}
Both hulls are $R$-modules, so $A$-linearity is equivalent to
$R$-linearity. Set $h=\operatorname{ht}(\mathfrak pR)$; then
$\dim R=h+1$. Choose parameters $x_1,\ldots,x_h$ in the regular
local ring $R_{\mathfrak p}$. In its top local cohomology,
which is its injective hull, the cyclic submodules generated by
\[
 e_n=[1/(x_1\cdots x_h)^n],\qquad
 J_n=(x_1^n,\ldots,x_h^n)R_{\mathfrak p},
\]
exhaust $E_{\mathfrak p}$, and $\operatorname{Ann}(e_n)=J_n$.
The ideals $J_n$ are cofinal with powers of $\mathfrak pR_{\mathfrak p}$.
For $h=0$ use $e_n=1$ and $J_n=0$ in the fraction field.
These are the usual Gorenstein and local-cohomology identifications
\cite[Lemmas 47.21.3, 47.21.7, and proof of 47.21.10]{StacksGorenstein}.
Clearing finitely many coefficients over $R_{\mathfrak p}$ shows
that every finite set in $E_{\mathfrak p}$ lies in
$R(e_n/u)$ for suitable $n$ and $u\in R\setminus\mathfrak pR$.

Let $I_n=J_n\cap R$. This ideal is $\mathfrak pR$-primary,
so $R/I_n$ has dimension and depth one. The finite-free-resolution
bound for a regular local ring gives
$\operatorname{pd}_R(R/I_n)\leq h$
\cite[Proposition 10.110.1]{StacksRegularResolution}.
Applying $\Hom_R(R/I_n,-)$ to the residual resolution leaves only
\[
 E_{\mathfrak p}[I_n]\xrightarrow{\delta_{\mathfrak p\mathfrak q}}
 E_{\mathfrak q}[I_n]
\]
in degrees $h,h+1$. The last map is surjective. Its kernel
$W_n=\Ext_R^h(R/I_n,R)$ is finite over $R$.

Given a homomorphism $f$ and inputs in $R(e_n/u)$, lift
$f(e_n/u)$ through this surjection. Since the Artinian complete
intersection $R_{\mathfrak p}/J_n$ is Gorenstein,
$E_{\mathfrak p}[I_n]$ is generated by $e_n/u$ over
$R_{\mathfrak p}$. Thus some $c\in R_{\mathfrak p}$ gives
$\delta_{\mathfrak p\mathfrak q}(ce_n/u)=f(e_n/u)$.
This proves density. Every nonzero $c\in T$ acts surjectively on
$E_{\mathfrak p}$: its Matlis dual is multiplication by $c$ on
the regular complete local domain $T$. Since the displayed
surjection shows $\delta_{\mathfrak p\mathfrak q}\ne0$, the
coefficient map is injective.

Identify $E_{\mathfrak p}[I_n]$ with $(R/I_n)_{\mathfrak p}$
using $e_n$. The finite submodule $W_n$ localizes to this whole
ring, since $E_{\mathfrak q}[I_n]$ vanishes at $\mathfrak p$.
It is therefore commensurable with $R/I_n$: clearing its finite
set of denominators, and then lifting $1$ after localization,
gives $a_n,b_n\in R\setminus\mathfrak pR$ with
\[
 a_n(R/I_n)\subseteq W_n\subseteq b_n^{-1}(R/I_n).
\]
Vanishing on $e_n/u$ means that $c\bmod J_nT$ belongs to
$uW_n$. As $u$ varies, these subgroups are cofinal with
$uR+J_nT$. The ideal $(u,I_n)$ is $\mathfrak m$-primary, and
$\mathfrak m^N$ contains an element outside $\mathfrak pR$.
Thus this family is cofinal with
$\mathfrak m^N R+J_nT$. Cofinality of the $J_n$ with the
$\mathfrak p$-adic powers proves the asserted basis.
Finally the Hom module, being a closed subgroup of a product of
discrete copies of $E_{\mathfrak q}$, is separated complete.
The injective dense map therefore identifies it with the separated
completion just described.
\end{proof}

\begin{proposition}[I5: dense generation in every dimension; Proved]\label{i:dense}
The graded subalgebra of $\End_A^*(K)$ generated by its diagonal
degree-zero algebra and $\theta$ is dense in each degree for the
finite topology.
\end{proposition}
\begin{proof}
First fix $\mathfrak p\subseteq\mathfrak q$. Induct on the height
gap to approximate every map $E_{\mathfrak p}\to E_{\mathfrak q}$
on a finite set by finite sums of saturated paths of residual
arrows with diagonal coefficients. Gap zero follows from
Proposition~\ref{i:zero}; gap one is Proposition~\ref{i:adjacent}.
For larger gap work over $R=A_{\mathfrak q}$, of dimension $d$,
and put the prescribed inputs in a cyclic $R$-submodule $Re$,
where $e=e_n/u$ is as in the preceding proof. Its annihilator $I$
is $\mathfrak pR$-primary. Hence $R/I$ has positive depth, and
\cite[Proposition 10.110.1]{StacksRegularResolution} gives
$\operatorname{pd}_R(R/I)\leq d-1$.
The last differential of $\Hom_R(R/I,K_R)$ is consequently
surjective:
\[
 \bigoplus_{\substack{\operatorname{ht}(\mathfrak rR)=d-1\\
                      \mathfrak p\subseteq\mathfrak r\subset\mathfrak q}}
 E_{\mathfrak r}[I]\longrightarrow E_{\mathfrak q}[I].
\]
Lift $f(e)$; its preimage has finitely many nonzero components
$y_{\mathfrak r}$. The maps $Re\to E_{\mathfrak r}$ sending
$e$ to $y_{\mathfrak r}$ extend to maps
$g_{\mathfrak r}:E_{\mathfrak p}\to E_{\mathfrak r}$ by
injectivity. Thus $\sum_{\mathfrak r}
\delta_{\mathfrak r\mathfrak q}g_{\mathfrak r}$ agrees with
$f$ on $Re$. Approximate each $g_{\mathfrak r}$ on the prescribed
finite inputs by induction, and compose with
$\delta_{\mathfrak r\mathfrak q}$.

For the whole complex, finitely many inputs involve only finitely
many source summands, and their images involve only finitely many
target summands. Apply the preceding construction to these finitely
many component maps, insert endpoint idempotents, and sum. The
result has exactly the prescribed values on the original inputs.
Every path lies in the algebra generated by the diagonal algebra
and $\theta$, proving density.
\end{proof}

\paragraph{Refined higher-dimensional gate.}
Dense generation does not prove a presentation of a completed
algebra. It remains to establish the quotient topology and uniform
continuity needed to extend the universal map from the dense path
subalgebra, and to identify the entire closed ideal of relations
with the ideal generated by $\theta^2$ and incidence.
Yekutieli's Theorem~1.14, p.~9 of the opened author version,
constructs an adelic DG action; Lemma~1.11, pp.~7--8, is an
approximation statement \cite{Yekutieli2003}. An action does not
itself identify its kernel or supply the required quotient topology.
The full presentation in dimension at least three remains
\emph{Undetermined}.
